\begin{afterword}
\summaryhead

The post returns to Aristotelian categories, defined by lists of properties. It invents labels for combinations of lettered properties: ``zawa'' is A, C and D; ``yokie'' is B and E; ``xippo'' is E but not D; and so on up to ``merlacdonian.'' Replacing each label by its definition shows that the labels carry no information of their own. The rules of the system work on the definitions alone.

Applied to the syllogism about Socrates, with ``human'' defined as a mortal featherless biped, the replacement reveals a tautology. The label hid the premise and made the conclusion look new. You cannot call Socrates human until you have seen him die.

The author traces the saying ``you can define a word any way you like'' to this Aristotelian picture and proposes a different proverb: ``Definitions don't need words.''

\responsehead

The post makes one point by an exercise, and the exercise works. The labelled statements and their expansions are consistent, and the reader sees a string of invented words turn into lists of letters. The syllogism then shows the same thing: with ``mortal'' in the definition of ``human,'' the conclusion was assumed in the premise.

That point about the syllogism is old. John Stuart Mill reported in 1843 that the ``adversaries of the syllogistic theory'' urged exactly this, with the same example: ``Socrates is mortal'' is presupposed in ``All men are mortal.'' The post does not claim the point as new, and ``The Parable of Hemlock,'' earlier in the book, made it at length.

Two claims go beyond the exercise. The first is historical: the idea that you can define a word any way you like ``came from the Aristotelian notion of categories.'' No source is given, and the idea is older. In Plato's \textsc{Cratylus}, Hermogenes holds that any name one gives is as correct as another. The post's reason, that an agent following the Aristotelian rules would never reach a contradiction, explains why the rules allow the idea. It does not show where the idea began.

The second is that labels are ``only there to create the illusion of inference.'' Earlier the post called a label ``a human convenience (or inconvenience).'' Its own examples show both sides: a label can hide a premise, and it can also say in one word what the definition says in a nested list. The exercise supports the weaker claim, that labels add no information.

The closing proverb does useful work. It separates two claims that the familiar saying runs together: that labels are dispensable, which the post accepts, and that definitions have no consequences, which it rejects.

In short: a clear exercise showing that labels add nothing to definitions, making an objection to the syllogism that Mill reported in 1843, with an unsourced history of the idea it attacks, and a claim that labels are ``only'' illusion that goes beyond what the exercise shows.
\end{afterword}
